% Thesis roadmap (Chapter 1): problems of a streaming C2ST, the component of EADD
% that addresses each, the hypothesis that tests it, and where it is answered.
\begin{tikzpicture}[
    font=\small,
    c/.style={rectangle, draw=black!65, rounded corners=2pt, align=center,
              text width=3.3cm, minimum height=1.0cm, inner sep=3pt, fill=white},
    h/.style={c, fill=gray!12, text width=2.6cm},
    arr/.style={-{Stealth[scale=0.9]}, semithick, black!70},
    hd/.style={font=\footnotesize\bfseries, text=black!75}
]
\node[hd] at (0,0.9) {Problem of D3};
\node[hd] at (4.4,0.9) {Component of EADD};
\node[hd] at (8.4,0.9) {Hypothesis};
\node[hd] at (11.6,0.9) {Answered in};
\foreach \i/\pr/\co/\hy/\ch in {
    0/{sliding reference absorbs slow drift}/{reservoir reference since the last alarm}/{H1: slow drift}/{4.3, 5.1},
    1/{fixed AUC threshold, false alarms on autocorrelated data}/{permutation test with the history null}/{H2: false alarms}/{4.4, 5.2},
    2/{a full permutation test is costly}/{adaptive sequential test (ASPT)}/{H3: cost}/{4.5, 5.3},
    3/{no information on which features drifted}/{TreeSHAP ranking at an alarm}/{H4: attribution}/{4.6, 5.4}} {
    \node[c] (p\i) at (0,-1.3*\i) {\footnotesize \pr};
    \node[c] (c\i) at (4.4,-1.3*\i) {\footnotesize \co};
    \node[h] (h\i) at (8.4,-1.3*\i) {\footnotesize \hy};
    \node[font=\footnotesize, text=black!75] (s\i) at (11.6,-1.3*\i) {Sections \ch};
    \draw[arr] (p\i) -- (c\i);
    \draw[arr] (c\i) -- (h\i);
    \draw[arr] (h\i) -- (s\i);
}
\end{tikzpicture}
